
\documentclass[letter, 11pt]{article}

    %%%%%%%%%%%%%%%%%%%%%%%%%%%%%%
    \oddsidemargin  0.0in
    \evensidemargin 0.0in
    \textwidth      6.8in
    \headheight     0.0in
    \topmargin      0.0in
    \textheight=9.0in
    \parindent=0in
    %%%%%%%%%%%%%%%%%%%%%%%%%%%%%%


\renewcommand{\familydefault}{\sfdefault} % Police 

\usepackage{amsthm,amsfonts,amssymb,amsmath, xcolor, enumitem}
\usepackage{graphicx}   %To include graphics
\usepackage[small,compact]{titlesec}

\usepackage[utf8]{inputenc}

\usepackage[french]{babel}

\usepackage{fancyhdr}
\pagestyle{fancy}

\newcommand{\var}{\mathbb{V}\text{ar}} %Symboles programmés pour simplifier le texte



\usepackage{version} % Pour faire un texte troué pour les étudiants
\includeversion{etudiant}\excludeversion{professeur} %version etudiant
%\includeversion{professeur}\excludeversion{etudiants} %version prof
\newcommand{\etudiant}[1]{\begin{etudiants}#1 \end{etudiants}}
\newcommand{\prof}[1]{\begin{professeur}\color{blue} {\bf Solution :} #1 \end{professeur}}


\fancyhead[L]{STT-2100 \\ Régression}
\fancyhead[R]{A.-S. Charest, E. Reny-Nolin, A.-A. Vallée \\ Université Laval}


\begin{document}

\vspace*{0.01 cm}
\begin{center}
{\Large Exercices - Chapitre 3}
\end{center}




\begin{enumerate}[leftmargin=*, wide, label=\bf Exercice \arabic*:]

%%%%%%%%%%%%
% Question 1
\item Montrez quelques résultats du cours.
\begin{enumerate}
    \item L'expression $S_{xy}$ peut s'écrire sous différentes formes. Montrez que ces trois expressions sont égales,
    $$\sum\limits_{i=1}^n \left(x_i-\overline x\right)\left(Y_i-\overline Y\right) = \sum\limits_{i=1}^n \left(x_i-\overline x\right)Y_i = \sum\limits_{i=1}^n x_i\,Y_i - n\overline x\,\overline Y.$$
    \item Montrez que $\sum\limits_{i=1}^n (Y_i-\widehat Y_i)(\widehat Y_i-\overline Y) = 0$.
    \item Montrez que $\text{Cov}(\overline Y,\widehat\beta_1)=0$.
    \end{enumerate}

\prof{\begin{enumerate}
\item $$\begin{array}{ll}
    \sum\limits_{i=1}^n (x_i-\overline x)(Y_i-\overline Y)
    &= \sum\limits_{i=1}^n (x_i-\overline x)\,Y_i-\sum\limits_{i=1}^n(x_i-\overline x)\,\overline Y\\[3mm]
    &= \sum\limits_{i=1}^n (x_i-\overline x)\,Y_i-\overline Y \sum\limits_{i=1}^n(x_i-\overline x)\,\\[3mm]
    &= \sum\limits_{i=1}^n (x_i-\overline x)\,Y_i-\overline Y \left[\sum\limits_{i=1}^n x_i-n\overline x\right]\,\\[3mm]
     &= \sum\limits_{i=1}^n (x_i-\overline x)Y_i \\[3mm]
     &= \sum\limits_{i=1}^n x_i\,Y_i-\overline x\sum\limits_{i=1}^n Y_i \\[3mm]
    &= \sum\limits_{i=1}^n x_i\,Y_i - n\overline x\,\overline Y\\
    \end{array}$$

\item $$\begin{array}{ll}
    \sum\limits_{i=1}^n (Y_i-\widehat Y_i)(\widehat Y_i-\overline Y)
    &= \sum\limits_{i=1}^n Y_i\widehat Y_i- \overline Y\sum\limits_{i=1}^n Y_i- \sum\limits_{i=1}^n \hat Y^2_i+ \overline Y\sum\limits_{i=1}^n \widehat Y_i\\[6mm]
    &= \sum\limits_{i=1}^n Y_i\widehat Y_i-  \sum\limits_{i=1}^n \hat Y^2_i\\[6mm]
    &= \sum\limits_{i=1}^n Y_i [\overline Y - \widehat\beta_1\overline x +\widehat\beta_1\,x_i]- \sum\limits_{i=1}^n [\overline Y - \widehat\beta_1\overline x +\widehat\beta_1\,x_i]^2\\[6mm]
    &= \sum\limits_{i=1}^n Y_i [\overline Y +\widehat\beta_1(x_i-\overline x) ]- \sum\limits_{i=1}^n [\overline Y +\widehat\beta_1(x_i-\overline x) ]^2\\[6mm]
    &= \overline Y \sum\limits_{i=1}^n Y_i  +\widehat\beta_1\sum\limits_{i=1}^n(x_i-\overline x)Y_i -\left[ \sum\limits_{i=1}^n \overline Y^2 +\widehat\beta_1^2\sum\limits_{i=1}^n(x_i-\overline x)^2 + 2\,\overline Y\,\widehat\beta_1\sum\limits_{i=1}^n(x_i-\overline x)\right]\\[6mm]
    &= n\overline Y^2 +\widehat\beta_1\,S_{xy}  -\left[n \overline Y^2 +\widehat\beta_1^2\,S_{xx} \right]\\[6mm]
    &=0  
    \end{array}$$
\item  $$\begin{array}{ll}
    \text{Cov}(\overline Y,\widehat\beta_1)
    &=  \text{Cov}\left(\sum\limits_{i=1}^n\dfrac{Y_i}{n}, \sum\limits_{j=1}^n \dfrac{(x_j-\overline x)Y_j}{S_{xx}}\right) \\[6mm]
    &= \dfrac{1}{nS_{xx}}\text{Cov}\left(\sum\limits_{i=1}^n Y_i ,\sum\limits_{j=1}^n (x_j-\overline x)Y_j\right) \\[6mm]
    &= \dfrac{1}{nS_{xx}}\sum\limits_{i=1}^n\sum\limits_{j=1}^n (x_j-\overline x)\text{Cov}\left( Y_i ,Y_j\right) \\[3mm]
    \end{array}$$

    La covariance de $Y_i$ et de $Y_j$ est nulle si $i\ne j$ à cause de l'indépendance des observations.
    
    Si $i=j$, alors $\text{Cov}( Y_i ,Y_j)=\text{Cov}( Y_i ,Y_i)=\var(Y_i)=\sigma^2$

$$\begin{array}{ll}
    \text{Cov}(\overline Y,\widehat\beta_1)
    &= \dfrac{1}{nS_{xx}}\sum\limits_{i=1}^n(x_i-\overline x)\text{Cov}\left( Y_i ,Y_i\right) \\[6mm]
    &= \dfrac{1}{nS_{xx}}\sum\limits_{i=1}^n(x_i-\overline x)\sigma^2\\[6mm]
    &=0.
    \end{array}$$
    \end{enumerate}}



%%%%%%%%%%%%
% Question 2
\item On s\rq{}intéresse à la concentration de quinine d\rq{}une solution. Il peut être difficile de mesurer cette concentration, et plus facile de prélever une mesure de fluorescence. Étudions la relation linéaire entre ces variables, nommée courbe de calibration.\footnote{Source: Statistics and Chemometrics for Analytical Chemistry  de James N. Miller et Jane C. Miller.}

\begin{table}[h!]
\centering
\begin{tabular}{l|rrrrrr|r|r}
&\multicolumn{6}{|c|}{Observations}&Moyenne&Variance\\\hline
Concentration (ng/ml) & 0 & 10 & 20 & 30 & 40 & 50&25,0&350,0\\\hline
Fluorescence  & 4 & 22 & 44 & 60 & 75 & 104&51,5&1312,7
\end{tabular}
\end{table}

\begin{enumerate}
\item  Écrivez le modèle de régression linéaire correspondant à l\rq{}estimation de la courbe de calibration.
\item  Calculez les estimations des paramètres $\beta_0,\; \beta_1, \; \sigma^2$.
\item  Testez l\rq{}hypothèse $H_0 : \beta_1 =  0$ contre $H_1 : \beta_1 \ne 0$.
\item   Calculez le coefficient de détermination et interprétez cette valeur selon le contexte.
\item   On observe une fluorescence de 47 pour une solution dont on ignore la concentration. Calculez un intervalle de niveau 95$\%$ pour la concentration de quinine de cette solution.
\item   Le responsable du laboratoire n\rq{}est pas satisfait des résultats: il trouve que l\rq{}intervalle calculé est trop grand. Comment pourrait-on réduire la taille de cet intervalle, tout en gardant un niveau de confiance de 95\%?\\
\end{enumerate}

\prof{ \begin{enumerate}[a) ]\itemsep3mm
\item  La variable indépendante est la fluorescence, notée $x_i$ et la variable dépendante est la concentration, notée $Y_i$. Le modèle de la régression linéaire est 

$$ Y_i = \beta_0 + \beta_1 x_i + \varepsilon_i \qquad\mbox{ où } \;\;\varepsilon_i \overset{iid}{\sim}N(0,\sigma^2). $$

\item On calcule d'abord
\begin{align*}
%\overline{x} &= (4+22+44+60+75+104)/6 = 51.5\\
%\overline{y} &= (0+10+20+30+40+50)/6 = 25 \\
S_{xx} &= (n-1)\,s^2_x = 5\times 1312,7=6563,5 \\
S_{yy} &= (n-1)\,s^2_y = 5\times 350=1750 \\
S_{xy} &= \sum x_iY_i-n\overline x\overline Y=(4\times0 + \ldots + 104\times 50) -6\times 51,5\times 25 = 3375.
\end{align*}
Donc,
\begin{align*}
\widehat{\beta}_1 &= 3375 / 6563,5 = 0,5142 \\
\widehat{\beta}_0 &= 25 - 51,5\times0,5142 = -1,4817.
\end{align*}
Pour calculer l'estimateur de $\sigma^2$, il nous faut les valeurs prédites:
\begin{align*}
\widehat{y}_1 &= -1.4817 + 4 \times 0.5142 = 0.5752 \\
\widehat{y}_2 &= -1.4817 + 22 \times 0.5142 = 9.8309 \\
\widehat{y}_3 &= -1.4817 + 44 \times 0.5142 = 21.1435 \\
\widehat{y}_4 &= -1.4817 + 60 \times 0.5142 = 29.3708 \\
\widehat{y}_5 &= -1.4817 + 75 \times 0.5142 = 37.0839 \\
\widehat{y}_6 &= -1.4817 + 104 \times 0.5142 = 51.9959.
\end{align*}
Donc,
$$ \widehat\sigma^2 = \dfrac{1}{n-2} \left[(0-0.5752)^2 + \ldots +(50-51.9959)^2 = \frac{14.55}{4} \right]= 3.6375. $$


\item On calcule d'abord
 $$ \widehat{\var(\widehat{\beta}_1)}  = \frac{\widehat\sigma^2}{S_{xx}} = \frac{3.6375}{6563.5} = 0.00055.$$
La statistique du test est
$$ t_{obs} = \frac{0.5142 - 0}{\sqrt{0.00055}} =  21.84.$$
On rejette $H_0$ si cette statistique est plus grande en valeur absolue  que $t_{n-2,1-\alpha/2} = t_{4,0.025} = 2.776$. On a que $t_{obs} = 21,84 > 2,776$ et on rejette l'hypothèse nulle. Le coefficient $\beta_1$ est significativement différent de 0.

\item
$$ R^2 = \dfrac{SC_{reg}}{SC_{tot}} = \dfrac{S^2_{xy}/S_{xx}}{S_{yy}}=\frac{(0.5752 -25)^2 + \ldots + (51.9959 - 25)^2}{S_{yy}} = \frac{1735.45}{1750} = 0.9917. $$
Donc, 99,17\% de la variation dans les valeurs de concentration est expliquée par la différence de fluorescence des solutions. 

\item On s'intéresse à une solution en particulier, et donc on veut un intervalle de prévision.% On utilise la formule (2.20) dans les notes et on obtient
\begin{align*}
 (-1.4817  + 0.5142\times 47) & \pm 2.776 * \sqrt{3.6375} \sqrt{ 1 + \frac{1}{6} + \frac{(51.5-47)^2}{6563.5} }\\
 22.6860 & \pm 5.7262.
\end{align*}
L'intervalle est   $(16.95, 28.41)$.


\item Il faudrait augmenter le nombre de solutions testées et utiliser des valeurs de fluorescence variées. 

\end{enumerate}
}



\end{enumerate}

\end{document}

